% !TeX root = ../../../main.tex
\subsection{Discussion and Conclusions}
We developed an analytical pipeline for processing free-text biomedical case reports into Sepsis-3 textual time series using large language models.
The resulting T2S2 corpus contains 2,139 open-access PMOA case reports represented as event--timestamp trajectories, providing a retrospective resource for temporal modeling of sepsis narratives.
Evaluation on the clinician-annotated \texttt{sepsis-40} subset shows strong performance, with GPT-5 achieving the strongest overall results and DeepSeek-R1 and Llama 3.3 providing competitive open-model baselines with different tradeoffs across event recovery, temporal ordering, and absolute timestamp accuracy. 
Additional analyses, including document-level bootstrapping, clinician--clinician agreement, sepsis onset inference, hallucination auditing, self-consistency evaluation, and cross-model agreement on \texttt{sepsis-100} and the full corpus, further contextualize the reliability and limitations of the extracted trajectories.
These results support T2S2 as a retrospective research corpus for downstream tasks such as risk forecasting, causal reasoning, and disease trajectory characterization, while making clear that it is not intended as a deployment-ready clinical extraction system.
% We developed an analytical pipeline for processing free-text biomedical corpora to construct Sepsis-3 textual time series using large language models, and demonstrated strong performance characteristics, including high match rates and concordance, as well as minimal temporal discrepancies.
% We plan to release the full cohort annotations upon acceptance, providing a novel Sepsis-3 cohort with rich textual information, annotated with timestamps, which may be useful for downstream tasks such as risk forecasting, causal reasoning, and disease trajectory characterization. 
% In PMOA samples \texttt{sepsis-40} and \texttt{sepsis-100}, performance levels were strong with respect to clinician annotations (\texttt{sepsis-40}) and inter-LLM agreement (\texttt{sepsis-100}).

% There are several limitations to our work. 
% First, our product, the textual time series corpus for sepsis, is a clinical research informatics cohort, and possesses its characteristic limitations.  
% It should not be considered a representative sample, since case report corpora possess biases, including but not limited to publication biases, rare finding over-representation, and revisionism. 
% Second, while we analyzed multiple data sources, we are limited by the number of clinician annotations that are challenging and expensive to obtain. 
% % Our hope is that a revision of strong first-pass LLM annotations will be more straightforward and a human-in-the-loop process will result in improved clinical time series datasets. 
% Third, the substantial use of LLMs in the analytic pipeline can make error analysis more difficult, as their errors can be subtle and inaccurate.  
% Furthermore, our error analysis (Appendix~\ref{apd:error-analysis}) attributes the systematic under-identification of events in i2m4 (Figure~\ref{fig:manual_vs_l33_event_count}) to the prevalence of semi-structured elements, such as tables and unexpanded abbreviations, within the discharge summaries. This indicates that while foundation models excel with standard narrative prose, extracting events from the dense, non-standard syntax often found in other clinical notes requires tailored prompting strategies.

There are several limitations to our work.
First, T2S2 is a clinical research informatics corpus constructed from published case reports and should not be considered a representative sample of all patients with sepsis.
Case report corpora are affected by publication bias, rare finding over-representation, selective reporting, and retrospective revisionism.
The corpus is therefore best suited for studying retrospective sepsis trajectories, developing temporal modeling methods, and analyzing narrative representations, rather than estimating population-level incidence or deployment-ready clinical performance.

Second, while \texttt{sepsis-40} provides a substantially expanded human gold standard, exhaustive clinician validation of the full 2,139-report corpus remains infeasible.
The additional validation analyses reduce but do not eliminate this limitation: hallucination auditing provides direct evidence of factual grounding in a sampled subset, self-consistency evaluates stability across repeated generations, and full-corpus cross-model agreement assesses reproducibility across model families.
However, these analyses do not replace clinician adjudication of every extracted event.
A complete characterization of corpus-wide errors would require a larger human-in-the-loop adjudication framework in which clinicians resolve model disagreements and temporal ambiguities.

Third, the substantial use of LLMs in the analytic pipeline can make error analysis difficult because errors may be subtle, clinically plausible, or concentrated in ambiguous temporal expressions.
Our hallucination audit suggests that unsupported event fabrication is uncommon in the audited subset under the extractive formulation, but timestamp errors remain possible even when the event itself is text-grounded.
This distinction is important: the primary failure mode in this task is often not hallucinating a clinical finding, but assigning the wrong time to a real finding, especially when timing requires chained relative reasoning or interpretation of vague duration phrases.

Fourth, the i2m4 analysis should be interpreted as a cross-schema evaluation.
The lower event match rates on i2m4 arise in part from mismatch between fragmented i2b2-style temporal concepts and T2S2-style standalone clinical findings.
For example, section headers, isolated abbreviations, or partial phrases can appear as i2b2 concepts, whereas the T2S2 representation intentionally favors clinically meaningful, compositional findings.
This means that i2m4 is useful as a legacy benchmark and stress test, but it is not a clean estimate of T2S2 extraction performance on clinical notes without re-annotation under the T2S2 schema.

While T2S2 provides a retrospective resource for studying sepsis trajectories, translating this pipeline into a prospective clinical tool would require a different extraction architecture.
Real-time clinical documentation is substantially noisier than finalized case reports or retrospective discharge summaries, and the available information changes continuously as care unfolds.
A prospective system would need to process note updates incrementally, integrate extracted textual findings with streaming structured EHR data such as vitals and laboratory results, calibrate predictions against prospective clinical outcomes, and evaluate safety under deployment-specific error modes.
Until such bridging work is completed, T2S2 should be viewed as a high-fidelity retrospective resource for temporal modeling and clinical trajectory analysis rather than as a clinical decision-support system.

Regarding generalizability to other medical conditions besides sepsis, our methodology could be extended in a straightforward manner.
Our motivation for choosing sepsis as the disease focus was to expand the analysis from the i2m4 dataset to a similar population in the PMOA corpus and maintain a well-defined phenotype.
Sepsis is a sensible choice as it is one of the most prevalent and important diagnoses in critical care medicine.  
However, there are several concerns when generalizing to other conditions.
For one, most other conditions are less acute in nature, meaning that the time distributions may be wider.
Also, the amount and type of information in such case reports may vary. 
Despite this, one benefit of AULTC compared to linear or squared discrepancies is that when intensive-care and chronic-disease cases are mixed, long-horizon discrepancies from chronic disease cases do not dominate the performance measure because of the log-time scaling.
Another generalizability concern is that for some conditions, less temporal information will be present.
This appears to be more acute in imaging studies, where the focus of the case is on the image content at a single snapshot in time.
For these types of studies, the longitudinal annotations become less meaningful compared to the image's information content.

Future work could focus on expanding this methodological framework.
While our current prompt augmentations did not yield systematic performance gains (Table \ref{ch06:tab:sens}), exploring advanced prompting paradigms alongside domain-specific fine-tuning could further refine the fidelity and expressivity of the extracted representations.
We anticipate that T2S2 will serve as a foundational resource for downstream healthcare applications, including causal inference, risk forecasting, and multimodal alignment.
The empirical utility of this representation is evidenced by our concurrent work, which leverages T2S2 from clinical narratives to accurately forecast future patient risk states \citep{noroozizadeh2026forecasting}.
Ultimately, advancing temporal AI in healthcare requires exploiting all available modalities; the extraction of structured temporal trajectories from narrative text provides a critical complement to existing structured data streams.

Subsequent iterations of this pipeline can be tailored to specific downstream objectives, each imposing distinct constraints.
For instance, causal modeling requires strict time-ordering, whereas forecasting demands precise absolute timestamps.
Furthermore, while our current formulation targets true physiological occurrence, an alternative paradigm could focus on the timing of clinical recognition.
Capturing these perceived occurrences could illuminate decision-making dynamics and facilitate research in dynamic resource allocation. 
Ultimately, integrating these narrative-derived timelines with structured physiological data offers a promising avenue to build more comprehensive, multimodal representations of patient trajectories.
% Ultimately, integrating these narrative-derived timelines with structured physiological data offers a promising avenue to reconcile true disease onset with clinical perception, advancing early sepsis detection.

% Subsequent iterations of this extraction pipeline can be optimized for specific downstream objectives, which impose distinct constraints on the generated time series. For example, time-ordering is a necessary condition for causality, whereas the precise absolute timestamp may be of high importance in forecasting and alignment. Finally, while our current formulation seeks to align clinical findings with their true physiological occurrence, an alternative modeling paradigm could focus on the timing of clinical recognition and documentation. Capturing these perceived occurrences could illuminate decision-making dynamics, facilitating research in multi-agent systems and dynamic resource allocation. Integrating these narrative-derived timelines with structured physiological data offers a promising avenue for reconciling true disease onset with clinical perception, ultimately advancing the early detection and management of sepsis trajectories.

